\chapter{绪论}
\label{chap:intro}

\section{研究背景}
\label{sec:intro:background}

电影是数字娱乐领域最具代表性的内容形态之一。近年来，流媒体平台、在线评分社区与开放电影数据库快速发展，电影相关数据在规模与维度上持续增长：用户评分与标签等行为数据能够反映受众偏好，影片类型、剧情简介与语言等元数据支撑内容理解，预算、票房与热度等商业指标则为市场分析提供重要依据。这些多源数据为电影检索、统计分析和智能推荐等应用提供了丰富素材。

然而，真实电影数据往往并不规整。不同数据源在字段命名、主键体系和缺失模式上存在显著差异，预算与票房等关键商业字段在公开数据集中常常不完整；用户评分矩阵又具有高度稀疏性，给协同过滤建模带来困难。与此同时，数据处理、模型训练与系统开发若彼此割裂，就难以形成从数据清洗到交互式服务落地的完整链路。因此，有必要围绕电影多源数据构建一套结构清晰、模块边界明确、可复现可扩展的分析与推荐系统。

基于上述背景，本文设计并实现 CineScope 系统。该系统以 MovieLens、TMDb 与 IMDb 数据为基础，贯通数据清洗、探索分析、建模预测、推荐服务与前端交互，旨在为电影数据分析提供一体化的工程化解决方案。

\section{问题定义}
\label{sec:intro:problem}

围绕电影数据分析与推荐这一应用场景，本文将研究问题概括为以下四个方面。

\textbf{（1）多源数据融合问题。} 系统需要从 MovieLens 获取用户行为数据，并结合 TMDb/IMDb 补充影片详情与财务元数据。如何在保证样本规模的前提下完成主键对齐、重复记录处理、缺失字段控制与质量审计，形成可供建模和在线服务使用的终版数据集，是首要研究问题。

\textbf{（2）票房建模问题。} 电影票房具有明显右偏分布，且同时受到预算、类型、评分统计、语言等多类特征影响。在有限样本条件下，如何选择合适的回归目标变换与模型结构，并通过可复现的离线训练流程获得稳定、可比较的评估结果，是本文关注的第二个问题。

\textbf{（3）推荐生成问题。} 实际应用中，用户既可能基于某部已看电影寻找相似影片，也可能希望根据个人历史评分获得个性化推荐。如何分别实现内容相似推荐与协同过滤推荐，并在工程上将推荐能力独立为可调用服务，同时处理冷启动与候选不足等场景，是推荐模块设计的核心问题。

\textbf{（4）系统交付问题。} 除算法模型外，电影数据分析能力还需要通过可交互的应用系统呈现。如何将电影检索、统计可视化、推荐代理、RAG 检索以及 Agent 交互整合为前后端分离、服务边界清晰的原型，是工程实现层面必须解决的问题。

\section{研究目标}
\label{sec:intro:objective}

围绕上述问题，本文以 CineScope 为研究对象，设定如下研究目标：

\begin{enumerate}
    \item \textbf{构建可复现的数据处理链路}：完成 MovieLens、TMDb 与 IMDb 数据的抽取、清洗与融合，产出包含电影主表、评分表、标签表、类型统计及质量报告在内的终版数据集，为后续分析与建模提供统一数据基础。
    \item \textbf{实现票房预测模型}：基于非零票房样本构建特征工程与训练流程，采用 XGBoost 与 LightGBM 集成方法完成票房回归建模，并在验证集上给出可比较的评估指标。
    \item \textbf{实现双通道推荐服务}：分别构建内容推荐与协同过滤推荐能力，通过独立微服务对外提供 HTTP 接口，支持相似电影召回与用户个性化推荐。
    \item \textbf{完成可交互系统原型}：基于 React 与 FastAPI 实现电影搜索、统计可视化、推荐展示及 RAG/Agent 扩展能力，验证“数据层—模型层—服务层—展示层”分层架构的工程可行性。
\end{enumerate}

通过实现上述目标，本文期望形成一套可运行、可验证、可扩展的电影数据分析系统。
